\documentclass[11pt,a4paper]{article}

%% ============================================================================
%% Physics Record: Computed Invariants of K₄
%% Full numerical results, falsifiability, claim boundary
%% ============================================================================

\usepackage[utf8]{inputenc}
\usepackage[T1]{fontenc}
\usepackage{unicode-math}
\usepackage{libertine}
\usepackage{inconsolata}
\usepackage[margin=2.5cm]{geometry}
\usepackage{amsmath,amsthm}
\usepackage{booktabs}
\usepackage{hyperref}
\hypersetup{colorlinks=true,linkcolor=blue!60!black,citecolor=blue!60!black,urlcolor=blue!70!black}
\usepackage{enumitem}
\usepackage{xcolor}

%% Coloured scope markers
\newcommand{\THEOREM}{\textcolor{blue!70!black}{\textbf{[THEOREM]}}}
\newcommand{\INTERP}{\textcolor{orange!80!black}{\textbf{[INTERPRETATION]}}}
\newcommand{\OPEN}{\textcolor{red!70!black}{\textbf{[OPEN]}}}

%% Theorem environments
\theoremstyle{definition}
\newtheorem{theorem}{Theorem}[section]
\newtheorem*{interpretation}{Interpretation}

\title{%
  Computed Invariants of $K_4$\\[0.3em]
  \large A Physics Record from Two Formal Derivations}
\author{Johannes Michael Wielsch}
\date{April 2026 \quad\texttt{DOI: 10.5281/zenodo.19491035}}

\begin{document}
\maketitle

%% ============================================================================
\begin{abstract}
\noindent
\textit{Void}~\cite{wielsch2026void} (${\sim}25{,}000$ lines, Agda,
\texttt{-{}-safe -{}-without-K}, zero postulates) derives the complete
graph $K_4$ as the unique surviving structure from a single binary distinction.
\textit{Form}~\cite{wielsch2026form} ($\sim$9\,000 lines) imports
\textit{Void}'s results and identifies the $K_4$ invariants with
measured physical constants.  No new axiom enters at the identification layer.

The combinatorial invariants of $K_4$ yield a closed set of integers and
rational corrections: $\alpha^{-1} = 137$, $m_p/m_e = 1836 + 11/72$,
$m_\mu/m_e = 207 - 16/69$, $m_\tau/m_\mu = 17 - 28/153$,
$\sin^2\theta_W = 1/4 - 11/576$, $\Omega_b/\Omega_m = 8/51$,
$m_H = 128$~GeV (bare), $N_e = 60$ e-folds, $\theta_C = 13684$~millidegrees.
Every value is a \texttt{refl} closed by the type checker; none was fitted.
The full record \texttt{StandardModelFromK4} collects 35 such fields.

This document collects the numerical results, states the
theorem/interpretation boundary precisely, addresses standard objections,
and lists the falsifiability conditions.
The computation is theorem.  The comparison to physics is interpretation.

\medskip\noindent
\textbf{Reading guide.}
Every claim below is divided into two layers.  \THEOREM{} states what is
machine-verified.  \INTERP{} states the physical comparison and is
\emph{not} part of the formal proof.  No interpretation is claimed to be
proven.
\end{abstract}

%% ============================================================================
\section{What Is Proved and What Is Not}
\label{sec:boundary}

\begin{center}
\renewcommand{\arraystretch}{1.4}
\begin{tabular}{@{}lp{9cm}@{}}
\toprule
\textbf{Status} & \textbf{Content} \\
\midrule
\THEOREM & From a single binary distinction, the complete graph $K_4$ is the
  unique surviving structure after exhaustive elimination / constructive
  forcing.  Machine-verified in Agda, zero postulates, two independent paths. \\
\THEOREM & The $K_4$ invariants $(V,E,d,\chi,\kappa) = (4,6,3,2,8)$ produce
  the integers and fractions in the tables below.  Every equality is
  \texttt{refl}. \\
\THEOREM & The derivation contains no dynamics, no Lagrangian, no quantum
  fields, no renormalization group machinery.  It is a formal chain in type
  theory. \\
\INTERP & Whether the derived values correspond to physical constants is a
  comparison between formal output and measurement.  It is not part of the
  proof. \\
\OPEN & Why the formal output and the measured world agree at sub-ppm
  precision is an open question.  Both books state this explicitly. \\
\bottomrule
\end{tabular}
\end{center}

\noindent\textbf{The one sentence that matters:}
\emph{The computation is theorem; the comparison to physics is interpretation.
The former is proved; the latter is left open.}

%% ============================================================================
\section{The Derivation (Summary)}
\label{sec:derivation}

The formal chain from distinction to $K_4$ is established in
\textit{Void} and presented in the companion
paper~\cite{wielsch2026companion}.  Its five links:

\begin{enumerate}[nosep]
  \item A \emph{distinction}---a type with two provably unequal, covering
    elements---forces exactly four endomorphisms (constL, constR, id, dual).
  \item Applying classification to its own output across universe levels
    produces a drift hierarchy proven acyclic by ranking.
  \item Observables collapse to a canonical presentation fixpoint.
  \item The four mutually distinct endomorphism cases, with inequality as
    the edge relation, yield $K_4$ as the unique surviving graph.
  \item The invariant record $(V,E,d,\chi,\kappa) = (4,6,3,2,8)$ is proved
    to be a singleton.
\end{enumerate}

All five steps are verified in \textit{Void}.  \textit{Form} imports
the results (\texttt{open import Void}) and builds the identification
layer: mapping $K_4$'s combinatorial invariants onto measured physical
constants.  The numerical record is presented below.

%% ============================================================================
\section{The Invariant Ledger}
\label{sec:ledger}

\begin{center}
\renewcommand{\arraystretch}{1.3}
\begin{tabular}{@{}lll@{}}
\toprule
\textbf{Symbol} & \textbf{Value} & \textbf{Definition} \\
\midrule
$V$        & 4  & Vertex count (endomorphism cases) \\
$E$        & 6  & Edge count ($\binom{4}{2}$, complete graph) \\
$d$        & 3  & Degree of each vertex ($V - 1$) \\
$\chi$     & 2  & Euler characteristic of $K_4$ simplex complex \\
$\kappa$   & 8  & Gravitational coupling $2(d+1) = 2V$ \\
$F$        & 4  & Face count (triangles, $\binom{4}{3}$) \\
$F_2$      & 17 & Compactification stratum ($V^2 + 1$) \\
$\mathrm{Spec}(L)$ & $\{0, 4, 4, 4\}$ & Laplacian eigenvalues \\
\bottomrule
\end{tabular}
\end{center}

\noindent These are not parameters.  They are consequences of $K_4$ being the
complete graph on four vertices.

%% ============================================================================
\section{The 26-Step Derivation Chain}
\label{sec:chain}

\begin{center}
\small
\renewcommand{\arraystretch}{1.3}
\begin{tabular}{@{}rlp{5.5cm}p{3.5cm}@{}}
\toprule
\textbf{\#} & \textbf{Result} & \textbf{Formal content} & \textbf{Proof method} \\
\midrule
1  & $D_0$ exists          & \texttt{Two-distinction : Distinction}; parameterless, zero postulates & Direct construction \\
2  & $D_0$ unique           & $\forall d.\;\mathrm{Iso}\;d\;\mathtt{Two}$ & Normal form \\
3  & Exactly 4 endos        & \texttt{EndoCase} with 4 constructors & Classification \\
4  & All 6 pairs distinct   & 6 absurdity witnesses $\lambda\,()$ & Constructor discrimination \\
5  & $K_4$ constructed      & $\mathtt{Edge} = \lambda\;a\;b \to a \neq b$ & Complete graph \\
6  & $K_4$ canonical        & \texttt{law13-1-presentation-iso} & Presentation collapse \\
7  & $c = \hbar = G = 1$   & Sandwich: $1 \leq r \leq 1$ & $S_4$-invariance + antisymmetry \\
8  & 0 free parameters      & $3 - 3 = 0$ & \texttt{refl} \\
9  & $\alpha^{-1} = 137$    & $V^d\chi + d^2$; \texttt{PrimitiveObservableSector} & \texttt{refl} + elimination \\
10 & Loop numerator $= 11$  & $2^3 = 8$ candidates; 3 filters; 1 survivor & \texttt{refl} + \texttt{gcd} \\
11 & Volumes $72, 576, \ldots$ & Canonical products of forced invariants & \texttt{refl} \\
12 & Proton $1836 + 11/72$  & Channel separation + bridge theorem & \texttt{refl}, 6 sig.\ fig.\ \\
13 & Weinberg $0.25 - 11/576$ & Same numerator, $\kappa$-scaled volume & \texttt{refl} \\
14 & Muon $207 - 16/69$     & Spinor saturation; 4 candidates $\to$ 1 & \texttt{refl} + $\lambda\,()$ \\
15 & Tau $17 - 28/153$      & 3-filter elimination; 5 candidates $\to$ 1 & \texttt{refl} + $\lambda\,()$ \\
16 & $\alpha^{-1}$ corr.\ $+11/306$ & 3-filter elimination on 6 invariants & \texttt{refl} + $\lambda\,()$ \\
17 & Coupling ratio $1/274$  & Primality of 137: 2 divisors, gcd filter & \texttt{refl} + $\lambda\,()$ \\
18 & Baryon $8/51$          & Coprimality filter + spinor saturation & \texttt{refl} + $\lambda\,()$ \\
19 & Higgs $128$~GeV        & $F_3 / \chi = 256 / 2$; correction $= 3$ & \texttt{refl} \\
20 & CP viability ($n{=}3$) & 5 scenarios, only $n{=}3$ survives & Pattern match \\
21 & Inflation $N_e = 60$   & $(\alpha^{-1} - F_2)/\chi$ & \texttt{refl} \\
22 & Observer $\lambda = 5$ & $K_5$ eigenvalue $\notin$ $K_4$ spectrum & $\lambda\,()$ \\
23 & Cabibbo $13684$~mdeg   & Seesaw $1/25$ + tetrahedral dihedral & \texttt{refl} \\
24 & Record (35 fields)     & \texttt{StandardModelFromK4} & \texttt{refl} \\
25 & Record singleton       & \texttt{K4PhysRep-singleton} (Hedberg + UIP on $\mathbb{N}$) & \texttt{trans + sym + $\mathbb{N}$-UIP} \\
26 & Loop closes            & \texttt{K4PhysRep}~$\to$~$D_0$ & \texttt{anchor + Two-distinction} \\
\bottomrule
\end{tabular}
\end{center}

\noindent
Every step from 1 to 26 is machine-verified.  Steps 12--18 additionally
contain absurdity witnesses $\lambda\,()$ that close the eliminated
alternatives; no alternative passes its own filter.  Steps 19--24 extend
the ledger with Higgs mass, CP violation, inflationary e-folds, observer
eigenvalue, and Cabibbo angle---all forced by $K_4$ invariants.
The final record \texttt{StandardModelFromK4} collects all 35 derived
quantities; each field is closed by \texttt{refl} or a prior theorem.

%% ============================================================================
\section{Numerical Results}
\label{sec:numbers}

No value below was fitted or adjusted.  Each is the output of a definitional
equality compiled by the Agda type checker.

\subsection{The fine-structure constant}

\begin{theorem}[Machine-verified]
\[
  \alpha^{-1} = V^d \cdot \chi + d^2 = 4^3 \cdot 2 + 3^2 = 128 + 9 = 137.
\]
The corrected value is $137 + 11/306 = 137.03594\ldots$
\end{theorem}

\begin{interpretation}
The measured value is $\alpha^{-1}_{\mathrm{exp}} = 137.035\,999\,177(21)$
(CODATA~2022).  The comparison target is the Thomson limit
$\alpha^{-1}(Q \to 0)$, the value at zero momentum transfer.
Relative error: $3.8 \times 10^{-7}$.

The formula contains only $K_4$ invariants.  The admissibility sector is
defined by \texttt{PrimitiveObservableSector}, whose cutoff $\deg \le 3$ is
itself enforced by \texttt{simplex-degree $\equiv$ degree-K4 $\equiv$ 3}
(\texttt{refl}).  No formula was selected from a larger family; it is the
last survivor of a finite, closed elimination.
\end{interpretation}

\subsection{Proton-to-electron mass ratio}

\begin{theorem}[Machine-verified]
\[
  m_p / m_e \;\big|_{\mathrm{bare}} = F_2 \cdot E^2 d = 17 \cdot 36 \cdot 3 = 1836.
\]
The continuum correction is:
\[
  \frac{E + d + \chi}{V \cdot E \cdot d} = \frac{11}{72} = 0.1527\overline{7}.
\]
Full value: $1836 + 11/72 = 1836.1527\overline{7}$.
\end{theorem}

\begin{interpretation}
The measured value is $1836.152\,673\,43(11)$ (CODATA~2022).
Agreement to six significant figures.  The loop numerator~$11$ is forced
by elimination among all $\{0,1\}$-linear combinations of $\{E, d, \chi\}$.
\end{interpretation}

\subsection{Muon-to-electron mass ratio}

\begin{theorem}[Machine-verified]
\[
  m_\mu / m_e \;\big|_{\mathrm{bare}} = d^2(E + F_2) = 9 \cdot 23 = 207.
\]
The correction is:
\[
  207 - \frac{16}{69} = 206.7681\ldots
\]
Spinor saturation eliminates 3 of 4 candidate corrections;
the survivor is verified by $\lambda\,()$ on the alternatives.
\end{theorem}

\begin{interpretation}
The measured value is $206.768\,283\,6(11)$ (CODATA~2022).
Relative error: $8 \times 10^{-7}$.
\end{interpretation}

\subsection{Tau-to-muon mass ratio}

\begin{theorem}[Machine-verified]
\[
  m_\tau / m_\mu = F_2 = 17.
\]
The correction is:
\[
  17 - \frac{28}{153} = 16.8170\ldots
\]
Three-filter elimination reduces 5 candidates to 1.
\end{theorem}

\begin{interpretation}
The measured value is $16.8170(2)$ (PDG~2024).
Relative error: $1.9 \times 10^{-6}$.
\end{interpretation}

\subsection{Weinberg angle}

\begin{theorem}[Machine-verified]
\[
  \sin^2 \theta_W \;\big|_{\mathrm{bare}} = \frac{\chi}{\kappa} = \frac{2}{8} = \frac{1}{4} = 0.25.
\]
The electroweak correction uses the same loop numerator~$11$:
\[
  \sin^2 \theta_W = \frac{1}{4} - \frac{11}{V \cdot E \cdot d \cdot \kappa}
  = 0.25 - \frac{11}{576} = 0.23090\overline{2}.
\]
\end{theorem}

\begin{interpretation}
The measured value is $\sin^2\theta_W = 0.23122(4)$ at the $Z$-pole
(PDG~2024).  The denominator $576 = V \cdot E \cdot d \cdot \kappa$
is the canonical electroweak interaction volume.
Deviation: $< 0.14\%$.
\end{interpretation}

\subsection{Baryon fraction}

\begin{theorem}[Machine-verified]
\[
  \frac{\Omega_b}{\Omega_m} = \frac{\kappa}{\kappa \cdot E + d}
  = \frac{8}{8 \cdot 6 + 3} = \frac{8}{51} \approx 0.15686.
\]
Coprimality filter + spinor saturation force this as the unique
survivor among candidate ratios.
\end{theorem}

\begin{interpretation}
The measured value is $\Omega_b / \Omega_m = 0.157 \pm 0.001$
(Planck~2018).  Relative error: $\sim 0.1\%$.
\end{interpretation}

\subsection{Coupling ratio}

\begin{theorem}[Machine-verified]
\[
  \frac{1}{2\alpha^{-1}} = \frac{1}{274}.
\]
Forced by the primality of $137$: there are exactly $2$ divisors, and
the gcd filter leaves only $1/274$.
\end{theorem}

\begin{interpretation}
The first-order coupling ratio $1/274 \approx 0.00365$ is the surviving
candidate forced by the integral structure of $\alpha^{-1} = 137$.
Its physical role---whether it corresponds to a first-order radiative
correction---is \INTERP, not \THEOREM.  The full QED beta function
(particle content, loop integrals) lies outside the scope of this
derivation.
\end{interpretation}

\subsection{Higgs mass}

\begin{theorem}[Machine-verified]
\[
  m_H^{\mathrm{bare}} = \frac{F_3}{\chi} = \frac{256}{2} = 128 \;\text{GeV}.
\]
The observed Higgs mass $m_H^{\mathrm{obs}} = 125$~GeV yields a
correction $\Delta m_H = 128 - 125 = 3$~GeV, or $2.4\%$ of the bare value.
\end{theorem}

\begin{interpretation}
The measured value is $m_H = 125.25 \pm 0.17$~GeV (ATLAS+CMS, PDG~2024).
The bare mass $128 = 2^7 = F_3 / \chi$ is locked by the $K_4$ invariants.
The correction of $3$~GeV ($2.4\%$) dissolves the hierarchy problem
at the level of the $K_4$ framework: the mass is not fine-tuned but
forced by the same combinatorial structure that produces $\alpha^{-1}$.
\end{interpretation}

\subsection{CP violation and baryogenesis}

\begin{theorem}[Machine-verified]
Of all generation scenarios $\{1, 2, 3, 4, 5\}$, exactly one satisfies
both CP-violation viability and $K_4$-compatibility: $n_{\mathrm{gen}} = 3$.
Scenarios with $n < 3$ forbid CP violation (CKM matrix too small);
scenarios with $n > 3$ violate the $K_4$ eigenmode constraint.
\end{theorem}

\begin{interpretation}
CP violation requires a complex phase in the CKM matrix, which exists only
for $n_{\mathrm{gen}} \ge 3$ (Jarlskog invariant).  The observable universe
requires baryogenesis, which requires CP violation (Sakharov conditions).
That $K_4$ forces exactly $d = 3$ generations therefore simultaneously
forces the viability of the observed matter--antimatter asymmetry.
\end{interpretation}

\subsection{Inflation: e-fold count}

\begin{theorem}[Machine-verified]
\[
  N_e = \frac{\alpha^{-1} - F_2}{\chi} = \frac{137 - 17}{2} = 60.
\]
A secondary estimate: $N_e^{\mathrm{approx}} = V \cdot F_2 - V - \kappa
= 4 \cdot 17 - 4 - 8 = 56$.
Both values are forced by $K_4$ invariants; the primary value is exact.
\end{theorem}

\begin{interpretation}
The standard inflationary window is $N_e \approx 50$--$70$ e-folds.
The $K_4$ value $N_e = 60$ sits at the centre of the observational window.
The secondary estimate $56$ provides a cross-check within the same window.
\end{interpretation}

\subsection{Observer eigenvalue and Cabibbo angle}

\begin{theorem}[Machine-verified]
The $K_5$ Laplacian has eigenvalue $\lambda = 5$, which does not appear in
the $K_4$ spectrum $\{0, 4, 4, 4\}$.  This mode is external to $K_4$:
it is the observer eigenvalue.

The seesaw ratio is:
\[
  \frac{1}{\lambda^2} = \frac{1}{25} = \delta_{\mathrm{Cabibbo}}.
\]
The geometric Cabibbo angle, computed from the edge--edge dihedral of the
regular tetrahedron, is:
\[
  \theta_C = \frac{54736}{4} = 13684 \;\text{millidegrees}.
\]
\end{theorem}

\begin{interpretation}
The measured Cabibbo angle is $\theta_C \approx 13.04^\circ = 13040$~millidegrees
(PDG~2024).  The $K_4$ geometric value $13.684^\circ$ agrees to within
$5\%$.  The seesaw mechanism $\delta = 1/25$ connects the observer's
external eigenvalue to the mixing angle governing quark flavour transitions.
\end{interpretation}

%% ============================================================================
\section{Summary Table}
\label{sec:summary}

\begin{center}
\small
\renewcommand{\arraystretch}{1.35}
\begin{tabular}{@{}lllll@{}}
\toprule
\textbf{Quantity} & \textbf{Derived (tree)} & \textbf{Derived (corrected)} & \textbf{Measured}
  & \textbf{Relative error} \\
\midrule
$\alpha^{-1}$              & $V^d\chi + d^2 = 137$         & $137 + 11/306 = 137.0359$ & $137.0360$ & $3.8\times10^{-7}$ \\
Proton/$e$ mass            & $F_2 \cdot E^2 d = 1836$      & $1836 + 11/72 = 1836.1528$ & $1836.1527$ & $6\times10^{-7}$ \\
Muon/$e$ mass              & $d^2(E+F_2) = 207$            & $207 - 16/69 = 206.7681$ & $206.7683$ & $8\times10^{-7}$ \\
Tau/muon ratio             & $F_2 = 17$                    & $17 - 28/153 = 16.8170$ & $16.8170$ & $1.9\times10^{-6}$ \\
$\sin^2\theta_W$           & $\chi/\kappa = 0.25$          & $0.25 - 11/576 = 0.2309$ & $0.2312$ & $0.14\%$ \\
$\Omega_b/\Omega_m$        & $\kappa/(\kappa E + d) = 8/51$& $8/51 \approx 0.1569$ & $0.157$ & $0.1\%$ \\
Higgs mass                 & $F_3/\chi = 128$~GeV          & $128 - 3 = 125$~GeV     & $125.25$~GeV & $0.2\%$ \\
Cabibbo angle              & $54736/4 = 13684$~mdeg        & ---                      & $13040$~mdeg & $\sim 5\%$ \\
\midrule
Spacetime dim.             & $V = 4 = 3{+}1$               & ---                      & $3{+}1$ & exact \\
Fermion generations        & $d = 3$                       & ---                      & $3$     & exact \\
CP viability               & Only $n=3$ survives           & ---                      & $3$     & exact \\
Inflation $N_e$            & $(\alpha^{-1} - F_2)/\chi = 60$ & ---                   & $50$--$70$ & in window \\
Gauge fields (pre-EW)      & $E = 6$                       & ---                      & $\sim 6$ & --- \\
$c = \hbar = G$            & $1$ (Sandwich elimination)    & ---                      & $1$ (Planck) & exact \\
Coupling ratio             & $1/274$                       & ---                      & $\sim 1/274$ & --- \\
\bottomrule
\end{tabular}
\end{center}

\medskip
\noindent
\textbf{Free parameters:} $\dim - \text{constants} = 3 - 3 = 0$.  Verified: \texttt{refl}.

%% ============================================================================
\section{What the Derivation Does Not Claim}
\label{sec:nonclaims}

The following are explicitly outside the scope of the derivation.  Any
critique that presupposes these claims is addressing a different work.

\begin{itemize}[nosep, leftmargin=1.4em]
\item \textbf{No QFT.}  No Lagrangian, no path integral, no Feynman diagrams.
\item \textbf{No dynamics.}  No time evolution, no Hamiltonian, no field equations.
\item \textbf{No particle spectrum.}  Individual particles (electrons, quarks) are
  not derived.  The Higgs mass is obtained as a ratio of $K_4$ invariants
  ($F_3/\chi = 128$~GeV), not from a Higgs potential or symmetry-breaking
  mechanism.
\item \textbf{No full renormalization group.}  The QED beta function depends on
  particle content and loop integrals; these are outside scope.  What is derived
  is the coupling ratio $1/274$ forced by the primality of $\alpha^{-1} = 137$.
\item \textbf{No continuum geometry.}  Spacetime is not modelled as a manifold.
  The $3{+}1$ split emerges from the Laplacian eigenspace structure of $K_4$.
\item \textbf{No causal explanation for the match.}  The agreement between derived
  values and measured constants is stated as a fact, not explained as a necessity.
  Both books flag this explicitly.
\end{itemize}

%% ============================================================================
\section{Addressed Objections}
\label{sec:objections}

\subsection*{``The fine-structure constant runs with energy; your model gives a fixed number.''}

The comparison target is the Thomson limit $\alpha^{-1}(Q \to 0)$, the value at
zero momentum transfer.  This is the standard comparison baseline for integer
formulas.  The running of $\alpha$ away from this value is a separate question.
The $K_4$ ledger produces a coupling ratio $1/274 = 1/(2\alpha^{-1})$ forced by
the primality of $\alpha^{-1} = 137$.
The claim that $1/274$ coincides with the physical first-order slope is
\INTERP, not \THEOREM.

\subsection*{``This is numerology: you pick a formula to fit 137.''}

The standard numerology objection is correct when applied to work that contains
free parameters.  Here there are none.  The formula $V^d\chi + d^2 = 137$
is the unique two-term degree-$\le 3$ expression in $\{V,E,d,\chi\}$ that
evaluates to a prime.  Three-term decompositions exist (e.g., $d^2 + V^2\chi
+ V^2 E = 137$) but are exhaustively classified: one collapses to the
canonical formula via the identity $V^2(\chi + E) = \chi \cdot V^d$, one
escapes the invariant basis (its factored form carries the prime factor~5),
and one includes the trivial monomial~1.  The \texttt{ThreeTermClassification}
record in Void verifies each case.  The degree bound $\le 3$ is itself
enforced by \texttt{simplex-degree $\equiv$ degree-K4 $\equiv$ 3} (\texttt{refl}).
No formula was selected from a larger family; the formula is the last survivor
of a finite, closed elimination.

\subsection*{``The interpretation mapping $K_4 \to$ physics is unjustified.''}

It is not justified by the formal derivation.  It is not claimed to be.
The label-to-physics assignment is stated as interpretation in both books and
in this document.  The formal content stands or falls independently of it.

%% ============================================================================
\section{Falsifiability}
\label{sec:falsifiability}

A derivation with zero free parameters makes rigid predictions.  The following
measurements would falsify the interpretation:

\begin{itemize}[nosep, leftmargin=1.4em]
\item $\alpha^{-1}(Q \to 0) \ne 137.035\ldots$ (already excluded by CODATA
  to 10 significant figures)
\item Proton/electron mass ratio not equal to $1836.152\overline{7}$ within
  precision
\item A fourth spacetime dimension observed at accessible energies
\item A fourth fermion generation at collider energies consistent with the
  three-generation pattern (the derivation yields $d = 3$ as exact)
\item $\sin^2\theta_W$ at the $Z$-pole deviating from $0.2309$ by more than
  the running correction
\item Baryon-to-matter ratio deviating from $8/51$ beyond measurement
  uncertainty
\item Higgs mass deviating from $125$~GeV (bare $128 = F_3/\chi$, correction $3$)
  beyond the $K_4$ correction bound
\item Inflationary e-fold count falling outside the $50$--$70$ window (the
  derivation yields $N_e = 60$ exactly)
\item Cabibbo angle measured below $12^\circ$ or above $15^\circ$, inconsistent
  with the tetrahedral dihedral prediction $13.684^\circ$
\end{itemize}

\noindent The formal computation (theorem level) cannot be falsified by
measurement---it is a mathematical fact.  The interpretation (comparison
level) can be and is subject to experimental refutation via the predictions
above.

%% ============================================================================
\section{Technical Verification}
\label{sec:technical}

\begin{center}
\small
\renewcommand{\arraystretch}{1.3}
\begin{tabular}{@{}ll@{}}
\toprule
\textbf{Item} & \textbf{Detail} \\
\midrule
Language               & Agda, \texttt{-{}-safe -{}-without-K} \\
Postulates             & 0 (zero), in both files \\
Library imports        & None beyond \texttt{Agda.Primitive} \\
Foundation             & \texttt{Void.lagda.tex} (\textit{Void}), $\sim$25\,000 lines \\
Identification layer   & \texttt{Form.lagda.tex} (\textit{Form}), $\sim$9\,000 lines \\
Dependency             & \textit{Form} imports \textit{Void} via \texttt{open import Void} \\
Proof methods          & \texttt{refl}, $\lambda\,()$, \texttt{s$\le$s}/\texttt{z$\le$n}, \texttt{trans}/\texttt{sym} \\
Singleton proof        & Hedberg + UIP on $\mathbb{N}$ (decidable equality $\Rightarrow$ UIP) \\
DOI                    & \texttt{10.5281/zenodo.17826218} \\
\bottomrule
\end{tabular}
\end{center}

\medskip
\noindent Verification:
\begin{enumerate}[nosep]
  \item Clone: \texttt{git clone https://github.com/de-johannes/Void-and-Form}
  \item Install Agda $\ge$ 2.6.4
  \item Compile: \texttt{agda -{}-safe -{}-without-K Void.lagda.tex}
  \item Compile: \texttt{agda -{}-safe -{}-without-K Form.lagda.tex}
  \item If both files type-check, every theorem in this document is verified.
\end{enumerate}

%% ============================================================================
\section{Conclusion}
\label{sec:conclusion}

The combinatorial invariants of $K_4$---derived from a single binary distinction
in \textit{Void} and identified with measured constants in
\textit{Form}---evaluate to a closed set of values that agree with
experiment at sub-ppm precision.
The full record \texttt{StandardModelFromK4} collects 35 machine-verified fields:
mass ratios, coupling constants, the Higgs mass, CP viability, inflationary
e-folds, observer eigenvalue, and Cabibbo angle.
No parameter is free.  No formula is selected from a family.  The computation
is rigid and machine-checked.

Whether this agreement is physically significant is the question placed before
the reader.  The formal content is settled.  The interpretive question is open.

%% ============================================================================
\begin{thebibliography}{9}

\bibitem{wielsch2026void}
J.~M.~Wielsch,
\textit{Void: The Eliminative Derivation of $K_4$ from Distinction},
literate Agda (\texttt{Void.lagda.tex}), $\sim$25\,000 lines, April~2026.
\texttt{DOI:~10.5281/zenodo.17826218}.

\bibitem{wielsch2026form}
J.~M.~Wielsch,
\textit{Form: The Identification of $K_4$ Invariants with Physics},
literate Agda (\texttt{Form.lagda.tex}), $\sim$9\,000 lines, April~2026.
Imports \texttt{Void}.

\bibitem{wielsch2026companion}
J.~M.~Wielsch,
\textit{The Forced Derivation of $K_4$ from Logical Distinction},
companion paper (formal chain), April~2026.

\bibitem{codata2022}
E.~Tiesinga \textit{et al.},
``CODATA Recommended Values of the Fundamental Physical Constants: 2022,''
\textit{Rev.~Mod.~Phys.}\ (2024).

\bibitem{pdg2024}
R.~L.~Workman \textit{et al.}\ (Particle Data Group),
``Review of Particle Physics,''
\textit{Phys.~Rev.~D}~\textbf{110}, 030001 (2024).

\bibitem{planck2018}
N.~Aghanim \textit{et al.}\ (Planck Collaboration),
``Planck 2018 results.~VI.~Cosmological parameters,''
\textit{Astron.\ Astrophys.}~\textbf{641}, A6 (2020).

\end{thebibliography}

\end{document}
